\documentclass[conference,a4paper]{IEEEtran}
% Six-page conference version. The full single-column version, with the related-work
% table, the parameter table and the latency figure, is STP_RAG_Research_Paper_v2_full.tex.

\usepackage[T1]{fontenc}
\usepackage[utf8]{inputenc}
\usepackage{microtype}
\usepackage{amsmath,amssymb}
\usepackage{booktabs,tabularx,array}
\usepackage{xcolor}
\usepackage{tikz}
\usetikzlibrary{arrows.meta,positioning,calc,fit}
\usepackage{pgfplots}
\usepgfplotslibrary{groupplots}
\pgfplotsset{compat=1.18}
\definecolor{stpfixed}{HTML}{333333}
\definecolor{stpsimilarity}{HTML}{999999}
\definecolor{stpvelocity}{HTML}{0072B2}
\definecolor{stpacceleration}{HTML}{E69F00}
\definecolor{stpfull}{HTML}{009E73}
\definecolor{stpcontext}{HTML}{CC79A7}
\definecolor{stpanswer}{HTML}{D55E00}
\pgfplotsset{
stpaxis/.style={
  width=0.44\linewidth,height=48mm,
  axis lines=left,axis line style={-},
  tick align=outside,tick label style={font=\footnotesize,/pgf/number format/fixed,/pgf/number format/precision=3},
  label style={font=\footnotesize},title style={font=\small\bfseries},
  grid=major,grid style={black!12},
  scaled ticks=false,
  legend style={draw=none,font=\footnotesize,column sep=5pt},
  legend cell align=left,
},
stpfixed/.style={color=stpfixed,mark=square*},
stpsimilarity/.style={color=stpsimilarity,mark=triangle*,mark options={rotate=180}},
stpvelocity/.style={color=stpvelocity,mark=*},
stpacceleration/.style={color=stpacceleration,mark=triangle*},
stpfull/.style={color=stpfull,mark=diamond*},
stpcontext/.style={color=stpcontext,mark=pentagon*},
}
\usepackage{algorithm}
\usepackage{algpseudocode}
\usepackage[hidelinks]{hyperref}
\usepackage[style=ieee,backend=biber,doi=false,url=false,isbn=false,maxbibnames=3,minbibnames=1,giveninits=true]{biblatex}
\addbibresource{references.bib}
\AtEveryBibitem{\clearfield{month}\clearfield{note}\clearlist{language}}
\renewcommand*{\bibfont}{\footnotesize}
\setlength{\bibitemsep}{0pt}

\DeclareMathOperator{\dist}{dist}
\DeclareMathOperator{\median}{median}
\DeclareMathOperator{\IQR}{IQR}
\DeclareMathOperator{\MAD}{MAD}
\DeclareMathOperator*{\argmax}{arg\,max}
\newcommand{\tok}{\operatorname{tok}}
\newcommand{\ctx}{\mathrm{ctx}}

\begin{document}

\title{STP-RAG: Semantic Transition Profile-Guided Adaptive Chunking and Selective Context for Retrieval-Augmented Generation}

\author{\IEEEauthorblockN{Mohammad Sameer, Rajeev Kumar Gupta, Mohit Dev Tomar, Dharm Raj, Parma Nand, Rani Astya}
\IEEEauthorblockA{Department of Computer Science and Engineering, Sharda University, Greater Noida, India\\
\{tech.geek.sameer, rajeevg0521, mohitdevtomar8, dharmraj4u, parmaastya, astyarani\}@gmail.com}}

\maketitle

\begin{abstract}
Retrieval-augmented generation (RAG) indexes documents as chunks, so the segmentation policy decides which statements can be retrieved together. We specify STP-RAG, which summarises local semantic change with a semantic transition profile (STP) of semantic velocity, positive acceleration, volatility and structural cues, and uses that one profile both to place chunk boundaries under length bounds and to select neighbouring context under a token budget. We also audit a six-arm pilot on HotpotQA and QASPER with a local phi3:mini generator. On HotpotQA, selective context raised full-STP answer F1 from 0.283 to 0.410; on QASPER, fixed-size chunking was best on every quality metric and context lowered F1 from 0.105 to 0.068. The audit shows why these figures are not confirmatory: 25 generated answers per dataset, 28.2\% of QASPER questions dropped, a velocity-only arm that kept structural weighting, and prompts of roughly 960 evidence tokens for fixed-size chunks against 345 for full STP and 945 with context. The HotpotQA gain is therefore consistent with restored prompt volume. We contribute a fully parameterised reference specification, a component-level reconciliation with the pilot code, and a pre-specified confirmatory protocol with matched token budgets, factorial ablations and paired, document-clustered inference, released with an open evaluation harness.
\end{abstract}

\begin{IEEEkeywords}
retrieval-augmented generation, text segmentation, adaptive chunking, context selection, evaluation methodology
\end{IEEEkeywords}

%==========================================================================
\section{Introduction}
\label{sec:intro}

Retrieval-augmented generation (RAG) conditions a language model on passages retrieved at answer time \cite{lewis2020rag}, and dense retrievers such as DPR made that retrieval effective for open-domain question answering \cite{karpukhin2020dpr}. Both take the passage as given. In practice every document must first be divided into chunks, and that division fixes whether a definition and the condition that limits it are retrieved together and how much text the generator must read to find one fact. Smaller chunks are focused but sever links between statements; larger chunks keep the links but dilute the evidence. Dense X Retrieval showed that retrieval and answer quality both depend on the retrieval unit \cite{chen2024densex}, and longer context windows do not remove the problem, because models use evidence less reliably when it is buried among irrelevant text \cite{liu2024lostmiddle}.

Prior work trains boundary detectors, changes the retrieval unit, or adds a representation layer. Less explored is the direct use of the \emph{dynamics} of local semantic change: how large a change is, whether it is accelerating, and how much it fluctuates. STP-RAG treats the cosine displacement between adjacent unit embeddings as semantic velocity, its positive change as acceleration and its local spread as volatility, and adds structural cues such as headings and paragraph breaks. We ask whether this one interpretable profile can steer two decisions that are usually designed separately: where to cut, and what neighbouring context to attach.

A pilot implementation was run before this specification was written. Our audit found its comparisons confounded by unequal prompt budgets, a contaminated baseline, small generation samples and dropped queries, so we report it as descriptive evidence and frame the research questions for a confirmatory run. \textbf{RQ1:} under a matched prompt budget, do acceleration, volatility, structure and selective context improve over velocity-only and fixed-size baselines? \textbf{RQ2:} how much of a quality difference is explained by chunk count and prompt tokens rather than by where boundaries fall? \textbf{RQ3:} is the effect of selective context stable across multi-hop Wikipedia questions and long scientific papers?

Our contributions are (i) a reference specification of STP-RAG with every parameter and numerical safeguard stated (Section~\ref{sec:method}); (ii) a component-level reconciliation of that specification with the audited pilot (Section~\ref{sec:audit}); (iii) the pilot aggregates with corrected labels and a prompt-token account (Section~\ref{sec:results}); and (iv) a pre-specified confirmatory protocol implemented in an open harness (Section~\ref{sec:protocol}).

\begin{figure*}[t]
\centering
\begin{tikzpicture}[
block/.style={draw,rounded corners=2pt,align=center,text width=24mm,minimum height=8.5mm,font=\scriptsize,fill=black!3},
stp/.style={block,fill=blue!7},
arr/.style={-{Latex[length=1.5mm]},semithick},
lab/.style={font=\scriptsize\bfseries},
x=33mm,y=-11mm]
\node[block] (units) at (0,0) {Sentence units,\\offsets, structure};
\node[stp]   (profile) at (1,0) {Transition profile\\$\mathbf{z}_i$ (Eqs.~1--4)};
\node[stp]   (seg) at (2,0) {Bounded\\segmentation (Alg.~1)};
\node[block,minimum height=19mm] (base) at (0,1.5) {Base chunks $c_i$\\with offsets\\and provenance};
\node[stp]   (ctxA) at (1,1) {Select context\\(Eqs.~7--8)};
\node[block] (embA) at (2,1) {Embed and index\\augmented $\tilde c_i$};
\node[block] (retA) at (3,1) {Retrieve to a\\token budget};
\node[block] (genA) at (4,1) {Generator\\prompt};
\node[block] (embB) at (1,2) {Embed and index\\base $c_i$};
\node[block] (retB) at (2,2) {Retrieve to a\\token budget};
\node[stp]   (ctxB) at (3,2) {Select context\\(Eqs.~7--8)};
\node[block] (genB) at (4,2) {Generator\\prompt};
\node[lab,above=0.5mm of genA.north east,anchor=south east] {Branch A};
\node[lab,below=0.5mm of genB.south east,anchor=north east] {Branch B};
\draw[arr] (units)--(profile); \draw[arr] (profile)--(seg);
\draw[arr] (seg.south) -- ++(0,-2mm) -| (base.north);
\draw[arr] (base.east |- ctxA) -- (ctxA);
\draw[arr] (base.east |- embB) -- (embB);
\draw[arr] (ctxA)--(embA); \draw[arr] (embA)--(retA); \draw[arr] (retA)--(genA);
\draw[arr] (embB)--(retB); \draw[arr] (retB)--(ctxB); \draw[arr] (ctxB)--(genB);
\draw[arr,dashed,blue!60!black] (profile.south) -- (ctxA.north);
\draw[arr,dashed,blue!60!black] (profile.north) -- (1,-0.55) -- node[above,font=\tiny]{profile reuse} (3.5,-0.55) -- (3.5,1.5) -- (3,1.5) -- (ctxB.north);
\end{tikzpicture}
\caption{STP-RAG with both context placements. Branch A changes what is indexed; branch B changes only the prompt. The pilot is consistent with branch B (Section~\ref{sec:audit}). Shaded blocks use the transition profile.}
\label{fig:architecture}
\end{figure*}

%==========================================================================
\section{Related Work}
\label{sec:related}

\textit{Retrieval and generation.} Adaptive-RAG, Self-RAG and CRAG adapt retrieval at query time through routing, self-reflection and correction \cite{jeong2024adaptiverag,asai2024selfrag,yan2024crag}, leaving documents unchanged. STP-RAG acts earlier, on how documents are cut and augmented, and is complementary to them.

\textit{Text segmentation.} TextTiling places boundaries at valleys of lexical cohesion \cite{hearst1997texttiling}, C99 clusters a ranked similarity matrix \cite{choi2000c99}, and TopicTiling scores boundaries from topic distributions \cite{riedl2012topictiling}. Supervised and transformer segmenters learn boundaries from labelled data \cite{koshorek2018textseg,lo2021transformer2}, and kernel change-point methods detect shifts in a representation space \cite{arlot2019kernelchangepoint}. $P_k$ and WindowDiff evaluate segmentations against a reference \cite{beeferman1999statistical,pevzner2002windowdiff}. Semantic velocity is closest to similarity-based segmentation; the STP adds second-order and dispersion terms, and none of these segmenters couples its boundaries to a retrieval context policy.

\textit{Retrieval units and context.} Dense X compares propositions, sentences and passages but fixes one unit type per index \cite{chen2024densex}. Late Chunking pools token embeddings of a long text into chunks so each embedding carries context, but selects no spans for the prompt \cite{gunther2024latechunking}. RAPTOR and GraphRAG add summary trees or entity graphs above the corpus at extra build cost \cite{sarthi2024raptor,edge2024graphrag}.

\textit{Evaluating RAG.} Knollmeyer et al.\ argue that retrieval and generation must be measured separately \cite{knollmeyer2026evaluation}, and Mereu et al.\ find that the benefit of retrieval varies by model and saturates with retrieval depth \cite{mereu2026specializing}. Both support reporting retrieval depth, context volume and per-stage metrics, which our protocol adopts. We evaluate on HotpotQA, multi-hop questions over Wikipedia paragraphs \cite{yang2018hotpotqa}, and QASPER, questions answered from full scientific papers \cite{dasigi2021qasper}.

%==========================================================================
\section{Reference Method}
\label{sec:method}

The specification is complete enough that an independent implementation produces the same chunks; Section~\ref{sec:audit} states what the pilot implemented.

\subsection{Semantic Transition Profile}

A document is a sequence of sentence units $D=(u_1,\ldots,u_n)$ with section, paragraph and character offsets; progression is measured in units, so $\Delta p_i=1$. Each unit is encoded by a frozen encoder and $L_2$-normalised to $e_i$. Semantic velocity, positive acceleration and volatility are
\begin{gather}
v_i = 1 - e_{i-1}^{\mathsf T}e_i \quad (i\ge 2), \label{eq:velocity}\\
a_i^+ = \max(0,\, v_i - v_{i-1}) \quad (i\ge 3),\ a_2^+=0, \label{eq:acceleration}\\
\sigma_i = \Big(\tfrac{1}{|W_i|}\textstyle\sum_{j\in W_i}(v_j-\bar v_{W_i})^2\Big)^{1/2}, \label{eq:volatility}
\end{gather}
where $W_i=\{\max(2,i-w+1),\ldots,i\}$ is a trailing window truncated at the document start. Clamping acceleration at zero matters: just after a velocity peak the signed change is strongly negative, and that drop is not evidence for a boundary. A binary vector $\mathbf{s}_i\in\{0,1\}^5$ marks whether $u_i$ starts a heading, paragraph, list item, display equation, or table or figure. Each continuous feature $x$ is scaled with development-split statistics,
\begin{equation}
\hat x_i = \frac{x_i-\median_{\mathrm{dev}}(x)}{\mathrm{scale}_{\mathrm{dev}}(x)},
\label{eq:normalise}
\end{equation}
where the scale is the IQR if it exceeds $\epsilon$, else $1.4826\,\MAD$, else the SD, else 1. The fallback matters because $a^+$ is zero at about half of all positions, so its IQR can collapse; implementations report how often each fallback fires. The profile is $\mathbf{z}_i=[\hat v_i,\widehat{a^+_i},\hat\sigma_i,\mathbf{s}_i]$.

\subsection{Bounded Adaptive Segmentation}

A boundary at $i$ is a cut between $u_{i-1}$ and $u_i$. Each candidate gets a score and probability
\begin{equation}
b_i = w_v\hat v_i + w_a\widehat{a^+_i} + w_\sigma\hat\sigma_i + \mathbf{w}_s^{\mathsf T}\mathbf{s}_i + \beta,\quad \pi_i = \frac{1}{1+e^{-b_i}}.
\label{eq:boundary}
\end{equation}
Let $\tok(s,j)$ count the tokens of $u_s,\ldots,u_{j-1}$. Algorithm~\ref{alg:segment} accepts a cut when $\pi_i>\tau_b$ and $\tok(s,i)\ge L_{\min}$. When adding $u_i$ would exceed $L_{\max}$, a cut is forced at
\begin{equation}
i^\star = \argmax_{j\in(s,i],\ \tok(s,j)\ge L_{\min}} \Big[\pi_j - \gamma\,\frac{|\tok(s,j)-L_t(s)|}{L_t(s)}\Big],
\label{eq:forced}
\end{equation}
with ties resolved to the later position, or at $i$ if no candidate is admissible. The target may adapt to volatility, $L_t(s)=\mathrm{clip}(L_t e^{-\eta\hat\sigma_s},L_{\min},L_{\max})$; $\eta=0$ gives a constant target. Every chunk keeps its unit identifiers and offsets.

\begin{algorithm}[t]
\caption{Reference STP segmentation}
\label{alg:segment}
\footnotesize
\begin{algorithmic}[1]
\Require tokens of $u_{1..n}$; $\pi_{2..n}$; $\tau_b,L_{\min},L_t,L_{\max},\gamma,\eta$
\State $C\gets[\,]$; $s\gets1$
\For{$i=2$ \textbf{to} $n$}
  \If{$\tok(s,i+1)>L_{\max}$} \Comment{$u_i$ overflows}
    \State $A\gets\{j\in(s,i]:\tok(s,j)\ge L_{\min}\}$
    \State $c\gets i$ \textbf{if} $A=\varnothing$ \textbf{else} maximiser of Eq.~\eqref{eq:forced}
    \State append $(s,c-1)$ to $C$; $s\gets c$
  \ElsIf{$\pi_i>\tau_b$ \textbf{and} $\tok(s,i)\ge L_{\min}$}
    \State append $(s,i-1)$ to $C$; $s\gets i$
  \EndIf
\EndFor
\State append $(s,n)$ to $C$; \Return $C$
\end{algorithmic}
\end{algorithm}

\subsection{Profile-Guided Selective Context}

For a base chunk $c_i$ spanning $[\mathrm{start}_i,\mathrm{end}_i]$, with embedding $\bar e_i$ the normalised mean of its unit embeddings, the profile sets how far the selector looks. A weak opening boundary suggests the chunk continues the preceding topic, so the horizon grows as the boundary weakens:
\begin{equation}
H_i^- = \lceil H_{\max}(1-\pi_{\mathrm{start}_i})\rceil,\quad H_i^+ = \lceil H_{\max}(1-\pi_{\mathrm{end}_i+1})\rceil.
\label{eq:horizon}
\end{equation}
Candidates within these horizons in the same document are scored with terms in $[0,1]$:
\begin{multline}
r(u_j,c_i) = \lambda_e\tfrac12(1+e_j^{\mathsf T}\bar e_i) + \lambda_r R(u_j,c_i) + \lambda_a A(u_j,c_i)\\
- \lambda_t T(u_j,c_i) - \lambda_d\,\dist(u_j,c_i)/H_{\max}.
\label{eq:contextscore}
\end{multline}
$R$ is the Jaccard overlap of entities, numbers and cross-references; $A$ is 1 when $u_j$ is in the same section; $\dist$ is the gap in units; and $T(u_j,c_i)=\max_k\pi_k$ over the boundaries between them is a \emph{transition barrier}, the second place the profile is reused: a candidate across a strong transition is penalised even if lexically similar. Setting $\lambda_t=0$ and $H_i^\pm=H_{\max}$ gives a relevance-only selector without STP information, which the protocol applies to every chunker as a control. Candidates with $r\ge\tau_r$ are added greedily within a per-chunk budget $B_{\ctx}$ and concatenated in source order.

Context can be placed in two ways, which must be declared (Fig.~\ref{fig:architecture}). In \emph{branch A} the augmented chunk is indexed, so context changes retrieval; in \emph{branch B} context is attached after retrieval and changes only the prompt. Retrieval credit is always computed on base text, and context tokens count against the prompt budget.

Proposed defaults, frozen in the protocol manifest before the confirmatory run, are $w=5$, $\epsilon=10^{-6}$, $\tau_b=0.5$, $(L_{\min},L_t,L_{\max})=(64,192,384)$ tokens, $\gamma=0.5$, $H_{\max}=6$, $(\lambda_e,\lambda_r,\lambda_a,\lambda_t,\lambda_d)=(1,0.5,0.25,1,0.25)$, $\tau_r=0.3$ and $B_{\ctx}=192$. Boundary weights are calibrated on the development split with Optuna \cite{akiba2019optuna} (200 trials, objective Recall@1024 tokens).


%==========================================================================
\section{Audit of the Pilot Implementation}
\label{sec:audit}

We audited the pilot repository in September 2026, covering its configuration, data splits, pipeline code and result notebook \cite{stpaudit2026}. Table~\ref{tab:reconcile} maps each component of the reference method to the pilot.

\begin{table}[t]
\caption{Reference specification versus the audited pilot. No pilot number validates a component not marked consistent.}
\label{tab:reconcile}
\centering
\footnotesize
\setlength{\tabcolsep}{2.5pt}
\begin{tabularx}{\columnwidth}{@{}>{\raggedright\arraybackslash}p{0.2\columnwidth}>{\raggedright\arraybackslash}X>{\raggedright\arraybackslash}p{0.2\columnwidth}@{}}
\toprule
Component & Pilot (audited) & Status \\
\midrule
Progression & Cumulative token positions, not $\Delta p_i=1$ & Differs \\
Encoder & bge-small in retrieval, bge-m3 in precomputation; unpinned & Conflict \\
Acceleration & Present; sign handling not confirmed & Unverified \\
Volatility & Present; window not released & Unverified \\
Normalisation & Stored statistics; no fallback & Partial \\
Structure & Scalar score, not binary vector & Differs \\
Velocity arm & Structural weight active & Contaminated \\
Context side & Predecessor units only & Differs \\
Profile in context & Horizon and barrier not released & Unverified \\
Placement & Retrieval unchanged by context & Branch B \\
Prompt evidence & Fixed top-3 chunks, not a token budget & Differs \\
Generator & phi3:mini; digest unrecorded & Partial \\
Records & Aggregate scalars only & Missing \\
\bottomrule
\end{tabularx}
\end{table}

\textit{Populations.} HotpotQA used 150 constructed ten-paragraph documents with 150 queries over one shared index. QASPER used 49 papers and 149 questions, of which 42 (28.2\%) were dropped because substring matching could not map their gold evidence; the 107 retained may favour short, extractive evidence.

\textit{Generation sample.} Answers were generated for 25 questions per dataset, so one exact answer moves EM by 0.04.

\textit{Prompt budget.} Every arm passed its top-3 chunks to the generator. Fixed-size chunks averaged about 300--350 tokens and full-STP chunks 95--121, so the generator received roughly 960, 345 and 945 evidence tokens for fixed-size, full STP and full STP with context. Answer quality is confounded with evidence volume.

\textit{Ablation design.} The arm labelled \texttt{velocity\_only} keeps the structural weight, so we relabel it velocity + structure. Arms are cumulative rather than factorial, and context was applied only to full STP.

\textit{Determinism.} Seeds 42, 123 and 999 left chunking, retrieval and generation unchanged, which invalidates a seed-level bootstrap \cite{efron1993bootstrap}. Small non-zero deviations appear only in F1, the metric most sensitive to one changed token; we read them as run-to-run generator noise, not uncertainty estimates. Confirming this needs per-seed predictions, which were not kept.

\textit{Models and records.} The retrieval configuration names bge-small-en-v1.5 \cite{xiao2023cpack} while precomputation defaults to bge-m3 \cite{chen2024bgem3}, and the phi3:mini \cite{abdin2024phi3} digest was not recorded. Splits, prompts and per-query predictions are not checksummed.

%==========================================================================
\section{Confirmatory Protocol}
\label{sec:protocol}

The protocol is fixed before any confirmatory result is seen; its manifest of parameters, splits, model digests and prompt template is committed with SHA-256 checksums, and any later change is reported as a deviation.

\textit{Populations.} HotpotQA uses 1,000 distractor-setting questions stratified by bridge and comparison type against one shared index; QASPER uses the full test split. Gold evidence is aligned by exact span lookup with a token-overlap fallback (Jaccard $\ge0.7$), so no answerable question is dropped. A chunk credits a gold span when at least half of its characters fall in the chunk's base text.

\textit{Budgets.} Retrieved chunks with any context are packed in rank order into $B_{\mathrm{prompt}}\in\{512,1024,2048\}$ tokens. We also report Recall@$B$ tokens and run fixed-size chunking at 100, 250 and 400 tokens and at full STP's mean chunk length.

\textit{Arms.} Table~\ref{tab:arms} lists the factorial segmentation arms. Each is crossed with no context, a naive equal-token window, relevance-only context and profile-guided context. Competitive baselines share the encoder, index and budgets: TextTiling, C99, a percentile semantic splitter, a recursive heading-and-paragraph splitter and a supervised segmenter \cite{hearst1997texttiling,choi2000c99,koshorek2018textseg}.

\begin{table}[t]
\caption{Factorial segmentation arms.}
\label{tab:arms}
\centering
\footnotesize
\setlength{\tabcolsep}{4.5pt}
\begin{tabular}{@{}lccccc@{}}
\toprule
Arm & $v$ & $a^+$ & $\sigma$ & $\mathbf{s}$ & $\eta\neq0$ \\
\midrule
Fixed-size (four sizes) & -- & -- & -- & -- & -- \\
Structure only & -- & -- & -- & \checkmark & -- \\
Velocity only & \checkmark & -- & -- & -- & -- \\
Velocity + structure (pilot arm) & \checkmark & -- & -- & \checkmark & -- \\
Acceleration only & -- & \checkmark & -- & -- & -- \\
Volatility only & -- & -- & \checkmark & -- & -- \\
Kinematics only & \checkmark & \checkmark & \checkmark & -- & -- \\
Full boundary, fixed target & \checkmark & \checkmark & \checkmark & \checkmark & -- \\
Full STP & \checkmark & \checkmark & \checkmark & \checkmark & \checkmark \\
\bottomrule
\end{tabular}
\end{table}

\textit{Metrics.} Recall@$K$ and MRR@10 use span credit, and EM and token F1 follow SQuAD normalisation \cite{rajpurkar2016squad}, with QASPER's official evaluator \cite{dasigi2021qasper}. Mechanism metrics are $P_k$ and WindowDiff against paragraph and section boundaries; the evidence split rate (gold spans spread over two or more chunks); context precision (share of attached tokens inside gold evidence); and F1 against prompt tokens.

\textit{Hypotheses.} Four primary hypotheses form one Holm-corrected family \cite{holm1979simple}. \textbf{H1:} at $B_{\mathrm{prompt}}=1024$, full STP with profile-guided context has higher F1 than length-matched fixed-size chunking with relevance-only context. \textbf{H2:} full STP has a lower evidence split rate than length-matched fixed-size chunking. \textbf{H3:} on the same base chunks, profile-guided context has higher context precision than relevance-only context. \textbf{H4:} kinematics only has higher Recall@1024 tokens than structure only.

\textit{Inference.} HotpotQA uses a paired bootstrap over queries (10,000 resamples); QASPER resamples whole papers \cite{field2007bootstrapping}. $p$-values come from paired sign-flip permutation tests \cite{dror2018hitchhiker}. If the per-query F1 difference has an SD near 0.35, 1,000 paired queries give a 95\% half-width of about 0.022. Every query writes a JSONL record of its retrieved chunks, context, prompt and metrics. The open harness implements this schema, budget packing, all metrics and tests, and reference versions of Algorithm~\ref{alg:segment} and Eq.~\eqref{eq:contextscore}; it has not been run on the pilot and contributes no results here.

%==========================================================================
\section{Pilot Results}
\label{sec:results}

\begin{figure*}[t]
\centering
\begin{tikzpicture}
\begin{groupplot}[stpaxis,group style={group size=2 by 1,horizontal sep=14mm},xmin=-0.095,xmax=0.145,ymin=-0.5,ymax=4.5,y dir=reverse,ytick={0,1,2,3,4},xtick={-0.05,0,0.05,0.10},xlabel={Difference from velocity + structure arm}]
\nextgroupplot[title={HotpotQA},yticklabels={Fixed,Similarity,Vel.+acc.,STP,STP+ctx},legend to name=stp-delta-legend,legend columns=2]
\addplot[black!65,no marks,forget plot] coordinates {(0,-0.5) (0,4.5)};
\addplot[xcomb,semithick,mark size=1.8pt,stpvelocity] coordinates {(-0.013,-0.14) (-0.020,0.86) (-0.007,1.86) (-0.013,2.86) (-0.013,3.86)};
\addlegendentry{Recall@5 (150 queries)}
\addplot[xcomb,semithick,mark size=1.8pt,color=stpanswer,mark=square*] coordinates {(0.014,0.14) (-0.020,1.14) (-0.061,2.14) (-0.061,3.14) (0.066,4.14)};
\addlegendentry{Answer F1 (25 queries)}
\nextgroupplot[title={QASPER},yticklabels=\empty]
\addplot[black!65,no marks,forget plot] coordinates {(0,-0.5) (0,4.5)};
\addplot[xcomb,semithick,mark size=1.8pt,stpvelocity] coordinates {(0.127,-0.14) (0.042,0.86) (0.005,1.86) (0.047,2.86) (0.047,3.86)};
\addplot[xcomb,semithick,mark size=1.8pt,color=stpanswer,mark=square*] coordinates {(0.027,0.14) (-0.029,1.14) (-0.023,2.14) (0.000,3.14) (-0.037,4.14)};
\end{groupplot}
\end{tikzpicture}
\par\smallskip
\pgfplotslegendfromname{stp-delta-legend}

\caption{Differences of rounded pilot means from the velocity + structure arm, for Recall@5 (150 HotpotQA and 107 QASPER queries) and answer F1 (25 queries per dataset). These are descriptive point estimates; no confidence intervals are implied.}
\label{fig:deltas}
\end{figure*}

Table~\ref{tab:pilot} reproduces the audited aggregates unchanged, with corrected labels. A deviation is shown only where it is non-zero; no intervals or tests are reported, because no per-query records exist.

\begin{table*}[t]
\caption{Pilot aggregates. EM gives exact answers out of 25 in parentheses; bold marks the best quality value per dataset. $^\dagger$Legacy \texttt{velocity\_only}, structural weight active. $^\ddagger$Legacy \texttt{velocity\_acceleration}, structural weight and sign unverified.}
\label{tab:pilot}
\centering\footnotesize
\setlength{\tabcolsep}{2pt}
\begin{tabular*}{\textwidth}{@{\extracolsep{\fill}}lccccccc@{}}
\toprule
Arm & Recall@5 & MRR & EM ($k$/25) & F1 & Chunks/doc. & Build (s) & Retrieval (ms) \\
\midrule
\multicolumn{8}{@{}l}{\textit{HotpotQA: 150 retrieval queries, 25 generated answers}} \\
Fixed-size & 0.987 & 0.917 & 0.200 (5) & 0.358 & 3.40 & 0.297{\scriptsize$\pm$0.178} & 1.717{\scriptsize$\pm$0.243} \\
Similarity threshold & 0.980 & 0.863 & 0.200 (5) & 0.324{\scriptsize$\pm$0.012} & 10.53 & 0.510{\scriptsize$\pm$0.006} & 3.972{\scriptsize$\pm$0.478} \\
Velocity + structure$^{\dagger}$ & \textbf{1.000} & \textbf{0.922} & 0.200 (5) & 0.344{\scriptsize$\pm$0.001} & 8.29 & 0.383{\scriptsize$\pm$0.003} & 2.564{\scriptsize$\pm$0.019} \\
Velocity + acceleration$^{\ddagger}$ & 0.993 & 0.912 & 0.120 (3) & 0.283{\scriptsize$\pm$0.002} & 9.04 & 0.439{\scriptsize$\pm$0.039} & 2.953{\scriptsize$\pm$0.111} \\
Full STP & 0.987 & 0.897 & 0.120 (3) & 0.283{\scriptsize$\pm$0.001} & 9.31 & 0.416{\scriptsize$\pm$0.004} & 2.822{\scriptsize$\pm$0.015} \\
Full STP + context & 0.987 & 0.897 & \textbf{0.240 (6)} & \textbf{0.410} & 9.31 & 1.064{\scriptsize$\pm$0.605} & 3.210{\scriptsize$\pm$0.411} \\
\midrule
\multicolumn{8}{@{}l}{\textit{QASPER: 107 aligned retrieval queries (42 of 149 dropped), 25 generated answers}} \\
Fixed-size & \textbf{0.379} & \textbf{0.226} & \textbf{0.040 (1)} & \textbf{0.132} & 12.43 & 0.390{\scriptsize$\pm$0.035} & 2.724{\scriptsize$\pm$0.174} \\
Similarity threshold & 0.294 & 0.214 & 0.000 (0) & 0.076{\scriptsize$\pm$0.001} & 29.53 & 0.795{\scriptsize$\pm$0.059} & 4.198{\scriptsize$\pm$0.589} \\
Velocity + structure$^{\dagger}$ & 0.252 & 0.187 & 0.000 (0) & 0.105{\scriptsize$\pm$0.001} & 41.80 & 1.005{\scriptsize$\pm$0.034} & 5.879{\scriptsize$\pm$0.223} \\
Velocity + acceleration$^{\ddagger}$ & 0.257 & 0.164 & 0.000 (0) & 0.082{\scriptsize$\pm$0.004} & 43.35 & 1.969{\scriptsize$\pm$1.405} & 5.088{\scriptsize$\pm$0.589} \\
Full STP & 0.299 & 0.180 & 0.000 (0) & 0.105{\scriptsize$\pm$0.002} & 47.96 & 1.335{\scriptsize$\pm$0.184} & 6.592{\scriptsize$\pm$1.072} \\
Full STP + context & 0.299 & 0.180 & 0.000 (0) & 0.068 & 47.96 & 1.527{\scriptsize$\pm$0.066} & 6.902{\scriptsize$\pm$0.796} \\
\bottomrule
\end{tabular*}

\end{table*}

\textit{RQ1: components.} On HotpotQA, velocity + structure has the best Recall@5 (1.000) and MRR (0.922), but it contains structure and cannot show whether velocity alone helps. Adding acceleration lowers F1 from 0.344 to 0.283 and EM from 5 to 3 of 25, and full STP matches that F1. Selective context raises full-STP F1 to 0.410 with Recall@5 unchanged, 14.5\% above fixed-size chunking; the EM difference behind this is three answers. On QASPER, fixed-size chunking is best on all four quality metrics: full STP reaches a Recall@5 of 0.299 against 0.379 ($-$21.1\%), and context lowers F1 from 0.105 to 0.068 ($-$35.2\%). Fig.~\ref{fig:deltas} plots each arm's difference from the velocity + structure arm.

\textit{RQ2: evidence volume.} Full STP makes 9.31 chunks per HotpotQA document against 3.40 for fixed-size, and 47.96 against 12.43 per QASPER paper. With a fixed number of retrieved chunks, finer segmentation delivers less evidence (Table~\ref{tab:budget}). The two highest-F1 HotpotQA arms delivered the most prompt tokens, and the context gain coincides with the prompt returning to the fixed-size volume. The most economical explanation is restored evidence volume, which H1 tests directly. Retrieval latency, 1.72--3.97\,ms on HotpotQA and 2.72--6.90\,ms on QASPER, grows with chunk count; on QASPER fixed-size chunking is both faster and more accurate. Timing conditions were not controlled.

\begin{table}[t]
\caption{Approximate evidence under fixed top-3 retrieval (audit estimates from mean chunk lengths) beside HotpotQA F1.}
\label{tab:budget}
\centering\footnotesize
\setlength{\tabcolsep}{4pt}
\begin{tabular}{@{}lccc@{}}
\toprule
Arm & Chunk (tok.) & Top-3 evidence & F1 \\
\midrule
Fixed-size & $\approx$300--350 & $\approx$960 & 0.358 \\
Full STP & $\approx$95--121 & $\approx$345 & 0.283 \\
Full STP + context & $\approx$115 + ctx & $\approx$945 & 0.410 \\
\bottomrule
\end{tabular}
\end{table}

\textit{RQ3: context reverses.} Full STP with and without context share base chunks, so their difference isolates the attachment step: $+$0.127 F1 on HotpotQA and $-$0.037 on QASPER. Without prompts or per-query records the pilot cannot say whether the reversal comes from coverage, noise, truncation or the generator, and the datasets differ in domain, length, question type and answer form. The confirmatory analysis breaks the effect down by evidence hops, evidence distance, cross-section evidence, document length and context length.

%==========================================================================
\section{Limitations and Conclusion}
\label{sec:limitations}

The pilot scored 25 answers per dataset, dropped 28.2\% of QASPER questions, has conflicting encoder settings and no per-query outputs, and its ablations are cumulative, so it supports no confirmatory or causal claim. Recall counted in chunks favours arms with larger chunks, a further reason for token-denominated metrics. The method also carries risks: calibrated weights may not transfer across domains, the profile depends on the encoder, and the transition barrier assumes strong transitions separate unrelated material, which may fail where closely linked topics alternate.

STP-RAG uses one profile of semantic velocity, positive acceleration, volatility and structure to drive both chunking and context selection, and every parameter and safeguard is stated. The audit shows that the pilot's HotpotQA gain coincides with restored prompt volume and a structure-weighted baseline, and that fixed-size chunking was more accurate and faster on QASPER. Comparing chunkers at a fixed number of retrieved chunks confounds segmentation with evidence volume. The confirmatory protocol removes that confound with matched budgets, full populations, pinned artifacts, factorial ablations, competitive segmenters and paired inference. The pilot aggregates, the scripts that generate every table, and the evaluation harness are released with this paper.

\printbibliography
\end{document}
